\documentclass[11pt]{article}

\usepackage[margin=1in]{geometry}
\usepackage{amsmath,amssymb,bm}
\usepackage{booktabs}
\usepackage[colorlinks=true,citecolor=blue,linkcolor=blue,urlcolor=blue]{hyperref}

\newcommand{\eps}{\epsilon_0({\bf k})}
\newcommand{\wls}{\omega_{\rm LS}({\bf k})}
\newcommand{\dd}{\delta_{\bf k}}
\newcommand{\Ck}{C_{\bf k}}
\newcommand{\Dk}{D_{\bf k}}
\newcommand{\Cnn}{C_1}

\title{Strong-coupling origin of the Luttinger-surface scale in the MMHK benchmark}
\author{Project note}
\date{June 15, 2026}

\begin{document}
\maketitle

\section{Conventions and numerical facts}

The extracted Luttinger-surface dispersion is fitted to the nearest-neighbor noninteracting dispersion
\begin{equation}
\eps=-2t(\cos k_x+\cos k_y),\qquad t=1,
\end{equation}
through
\begin{equation}
\wls\simeq a\eps+b .
\end{equation}
The present eta=0.05 data give the following asymptotic-looking scales:
\begin{align}
  a_{U\text{-only}}(U=40,V=0)&=-1.39286,\\
  a_{V\text{-only}}(U=0,V=12)&=-1.14265.
\end{align}
The purpose of this note is to show why these numbers are close to
\begin{equation}
  a_U\to 4\langle{\bf S}_i\cdot{\bf S}_j\rangle_{\rm spin},
  \qquad
  a_V\to -\frac87 .
\end{equation}
For the periodic $4\times4$ spin-$1/2$ Heisenberg cluster used by the MMHK block,
exact diagonalization gives
\begin{equation}
  E_0^{\rm Heis}=-11.2284832084,
  \qquad
  \Cnn=\frac{E_0^{\rm Heis}}{32}=-0.3508901003,
\end{equation}
where $32$ is the number of nearest-neighbor bonds.  Hence
\begin{equation}
  4\Cnn=-1.4035604011,
\end{equation}
which is already very close to the large-$U$ Luttinger-surface scale.

\section{Common two-pole/moment structure}

Inside the Mott or charge gap, the Green function zero is controlled by the cancellation
between particle-removal and particle-addition contributions.  A useful first-moment form is
\begin{equation}
  G({\bf k},\omega)
  \simeq
  \frac{w_h({\bf k})}{\omega-E_h({\bf k})}
  +
  \frac{w_p({\bf k})}{\omega-E_p({\bf k})},
\end{equation}
with $w_h+w_p=1$ for one spin species.  The zero is therefore
\begin{equation}
  \omega_0({\bf k})=w_h({\bf k})E_p({\bf k})+w_p({\bf k})E_h({\bf k}).
\end{equation}
Writing
\begin{equation}
  w_h=\frac12+\dd,
  \qquad
  w_p=\frac12-\dd,
\end{equation}
and
\begin{equation}
  \Ck=\frac{E_p+E_h}{2},
  \qquad
  \Dk=E_p-E_h,
\end{equation}
one obtains the exact two-pole identity
\begin{equation}
  \boxed{\omega_0({\bf k})=\Ck+\dd\Dk.}
  \label{eq:common-zero}
\end{equation}
The two asymptotic scales below differ only in the strong-coupling evaluation of
$\dd$, $\Dk$, and $\Ck$.

\section{$U$-only limit: spin-background scale}

At half filling and $U\gg t$, the low-energy sector of the Hubbard model maps to the
nearest-neighbor Heisenberg model with exchange $J=4t^2/U$.  The nearest-neighbor spin
correlator is
\begin{equation}
  \Cnn=\langle {\bf S}_i\cdot{\bf S}_{i+\hat x}\rangle.
\end{equation}
A Schrieffer--Wolff expansion gives the nearest-neighbor one-body density matrix, for a fixed
spin $\sigma$,
\begin{equation}
  \langle c^\dagger_{i\sigma}c_{i+\hat x,\sigma}\rangle
  =\frac{2t}{U}\left(\frac14-\Cnn\right)+O(t^3/U^3).
\end{equation}
Fourier transforming this nearest-neighbor expectation value gives
\begin{equation}
  n_{{\bf k}\sigma}=\frac12+\frac{2\Cnn-1/2}{U}\eps+O(t^3/U^3),
\end{equation}
so that
\begin{equation}
  \dd=n_{{\bf k}\sigma}-\frac12
  =\frac{2\Cnn-1/2}{U}\eps+O(t^3/U^3).
  \label{eq:delta-u}
\end{equation}
The addition--removal separation is $\Dk=U+O(t)$, so the weight-asymmetry contribution to
Eq.~\eqref{eq:common-zero} is
\begin{equation}
  \dd\Dk=(2\Cnn-1/2)\eps+O(t^2/U).
  \label{eq:deltad-u}
\end{equation}
The center of the two first moments is not zero.  Evaluating the first spectral moments, or
equivalently the $a_0$ Lanczos coefficients of the particle and hole continued fractions,
gives
\begin{equation}
  \Ck=(1/2+2\Cnn)\eps+O(t^2/U).
  \label{eq:center-u}
\end{equation}
Combining Eqs.~\eqref{eq:common-zero}, \eqref{eq:deltad-u}, and \eqref{eq:center-u},
\begin{equation}
  \boxed{
  \wls^{(U)}
  =\left[(1/2+2\Cnn)+(2\Cnn-1/2)\right]\eps
  =4\Cnn\eps+O(t^2/U).
  }
\end{equation}
For the present $4\times4$ MMHK spin background,
\begin{equation}
  \boxed{\wls^{(U)}\simeq -1.40356\,\eps,}
\end{equation}
consistent with the $U=40$ extracted scale $-1.39286$.

\begin{table}[h]
\centering
\caption{Numerical decomposition from the raw Lanczos files.  The first-moment zero uses the
$a_0$ coefficients and weights only.}
\begin{tabular}{cccc}
\toprule
$U$ & actual root scale & first-moment zero scale & $4\Cnn$ \\
\midrule
8  & $-1.33054$ & $-1.43530$ & $-1.40356$ \\
12 & $-1.30292$ & $-1.41479$ & $-1.40356$ \\
20 & $-1.36254$ & $-1.40559$ & $-1.40356$ \\
40 & $-1.39286$ & $-1.40368$ & $-1.40356$ \\
\bottomrule
\end{tabular}
\end{table}

\section{$V$-only limit: charge-defect scale}

For $U=0$ and $V\gg t$, the half-filled ground state is a checkerboard charge state: one
sublattice is doubly occupied and the other is empty.  A nearest-neighbor hop from a doubly
occupied site to an empty site creates a charge defect.  The defect energy denominator is
\begin{equation}
  \Delta_{\rm hop}=3(2V)+1V=7V.
\end{equation}
The three $2V$ terms are the interactions of the moved electron with the other occupied
nearest neighbors of the target site, and the remaining $V$ is the bond back to the source.
Thus the nearest-neighbor one-body density matrix has the leading amplitude
\begin{equation}
  \langle c^\dagger_{i\sigma}c_{i+\hat x,\sigma}\rangle
  =\frac{t}{7V}+O(t^3/V^3),
\end{equation}
and hence
\begin{equation}
  \dd=n_{{\bf k}\sigma}-\frac12
  =-\frac{1}{7V}\eps+O(t^3/V^3).
  \label{eq:delta-v}
\end{equation}
The particle-hole symmetric chemical potential is $\mu_{\rm ph}=4V$.  The dominant addition
and removal poles therefore lie near $\pm4V$, so
\begin{equation}
  \Dk=8V+O(t),
  \qquad
  \Ck=O(t^2/V).
  \label{eq:gap-v}
\end{equation}
Using Eq.~\eqref{eq:common-zero},
\begin{equation}
  \boxed{
  \wls^{(V)}
  =\dd\Dk+\Ck
  =-\frac{1}{7V}\eps\,(8V)+O(t^2/V)
  =-\frac87\eps+O(t^2/V).
  }
\end{equation}
This agrees with the extracted $U=0$ scales:
\begin{equation}
  a(V=4)=-1.14105,
  \qquad
  a(V=8)=-1.14240,
  \qquad
  a(V=12)=-1.14265,
\end{equation}
which converge to $-8/7=-1.142857\ldots$.

\section{Thermodynamic and larger-cluster estimates}

\subsection{Spin background for the $U$-only limit}

The $4\times4$ value $4\Cnn=-1.40356$ is a finite-block value.  In the thermodynamic
square-lattice spin-$1/2$ Heisenberg antiferromagnet, high-precision stochastic-series-expansion
QMC gives an energy density
\begin{equation}
  e_0=-0.669441857(7) \quad (J=1)
\end{equation}
for the Hamiltonian $H=J\sum_{\langle ij\rangle}{\bf S}_i\cdot{\bf S}_j$~\cite{sandvik2026highprecision}.
Since the square lattice has two bonds per site,
\begin{equation}
  \Cnn^{\infty}=e_0/2=-0.3347209285(35),
  \qquad
  4\Cnn^{\infty}=2e_0=-1.338883714(14).
\end{equation}
Thus a thermodynamic-limit spin background would predict a smaller large-$U$ scale,
approximately $-1.339$, whereas the present MMHK $4\times4$ block naturally gives the stronger
finite-cluster value $-1.404$.  This distinction is important: the observed $-1.4$ is not a
universal thermodynamic Hubbard number, but rather the nearest-neighbor spin correlation of the
finite MMHK block entering the strong-coupling zero.

\subsection{Charge background for the $V$-only limit}

There is an analogous thermodynamic estimate for the $V$-only result, but it is simpler than the
spin case because the $t=0$ charge background is classical.  On any even square lattice that
supports checkerboard order, the strong-coupling ground state has
\begin{equation}
  n_A=2,\qquad n_B=0,
\end{equation}
up to the two choices of occupied sublattice.  The leading hopping process is local: one electron
moves from a doubly occupied $A$ site to a neighboring empty $B$ site.  For square-lattice
coordination $z=4$, the defect denominator is
\begin{equation}
  \Delta_{\rm hop}=(2z-1)V=7V.
\end{equation}
The particle-hole separation is instead fixed by the particle-hole symmetric chemical potential,
\begin{equation}
  \mu_{\rm ph}=zV,
  \qquad
  \Dk\simeq 2zV=8V.
\end{equation}
Thus, for a general hypercubic bipartite lattice with coordination $z$, the same derivation gives
\begin{equation}
  \dd\simeq -\frac{1}{(2z-1)V}\eps,
  \qquad
  \wls^{(V)}\simeq -\frac{2z}{2z-1}\eps.
\end{equation}
For the square lattice,
\begin{equation}
  \boxed{a_V^{\infty}=-\frac{8}{7}.}
\end{equation}
Unlike the $U$-only coefficient, this leading $V$-only coefficient is not controlled by a
nontrivial finite-size spin correlation.  It is a local charge-defect denominator and therefore is
the same for the $4\times4$ block and for the thermodynamic checkerboard state, up to higher-order
virtual hopping corrections $O(t^2/V^2)$ and possible boundary effects on clusters that do not
support perfect checkerboard order.  The observed convergence
\begin{equation}
  -1.14105\;(V=4),\quad -1.14240\;(V=8),\quad -1.14265\;(V=12)
  \longrightarrow -1.142857\ldots
\end{equation}
is exactly the expected approach to this local thermodynamic value.

\begin{thebibliography}{9}
\bibitem{sandvik2026highprecision}
A.~W. Sandvik, ``High-precision ground state parameters of the two-dimensional spin-1/2
Heisenberg model on the square lattice,'' arXiv:2601.20189 (2026).
\end{thebibliography}

\end{document}
